\subsection{\texorpdfstring{第~\ref{sec:causality-in-lorentzian-manifolds}~节}{因果性}}\label{subsec:solutions-causality}

\begin{description}[font=\normalfont,leftmargin=2em,labelwidth=1.5em]
\RGsolution{ex:properties-of-chronological-and-causal-futures}{properties-of-chronological-and-causal-futures}
\begin{enumerate}[label=(\alph*),ref=\alph*,leftmargin=*]
\item 任取 $q\in I^+(p)$. 存在一条由 $p$ 连到 $q$ 的指向未来的类时曲线 $c:[0,1]\to M$. 取如 Proposition~\ref{prop:local-causal-structure-in-convex-normal-neighborhood} 中那样的 $q$ 的一个测地凸邻域 $V$, 再取 $s\in(0,1)$ 使 $c([s,1])\subset V$. 点 $r:=c(s)$ 相对于时空 $(V,g)$ 的时序未来 $I^+(r,V)$ 显然是开集 (一个开集在指数映射下的像), 且 $q\in I^+(r,V)\subset I^+(p)$. 因此 $q$ 是 $I^+(p)$ 的内点. 由于 $q$ 是任取的, $I^+(p)$ 为开集.
\item 令 $M$ 为从 $\mathbb{R}^2$ 中删去点 $(1,1)$ 后所得的空间. 若 $p$ 为原点, 则 $I^+(p)$ 是 Minkowski 锥
\[
  I^+(p)=\{(t,x)\in\mathbb{R}^2\mid t>|x|\}.
\]
但是, 无法沿因果曲线从 $p$ 到达直线 $t=x$ 上满足 $t>1$ 的任何一点. 因此,
\[
  J^+(p)=I^+(p)\cup\{(t,-t)\in\mathbb{R}^2\mid t\geq0\}
       \cup\{(t,t)\in\mathbb{R}^2\mid 0<t<1\}
\]
不是闭集.
\item 设 $q\in J^+(p)$. 则存在一条由 $p$ 连到 $q$ 的指向未来的因果曲线 $c:[0,1]\to M$. 设 $T$ 是沿 $c$ 平行的指向未来的类时向量场, 并取光滑映射 $\gamma:[0,1]\times(-\varepsilon,\varepsilon)\to M$, 使 $\gamma(0,t)=c(t)$ 且
\[
  \frac{\partial\gamma}{\partial s}(0,t)=tT_{c(t)}.
\]
令
\[
  \gamma'=\frac{\partial\gamma}{\partial s}=\gamma_*\frac{\partial}{\partial s}
  \qquad\text{且}\qquad
  \dot\gamma=\frac{\partial\gamma}{\partial t}=\gamma_*\frac{\partial}{\partial t},
\]
有
\[
  \nabla_{\gamma'}\dot\gamma-\nabla_{\dot\gamma}\gamma'
     =\gamma_*\left[\frac{\partial\gamma}{\partial s},\frac{\partial\gamma}{\partial t}\right]=0.
\]
因此
\begin{align*}
  \left.\frac{\partial}{\partial s}\right|_{s=0}\langle\dot\gamma,\dot\gamma\rangle
    &=\left.2\langle\nabla_{\gamma'}\dot\gamma,\dot\gamma\rangle\right|_{s=0}
     =\left.2\langle\nabla_{\dot\gamma}\gamma',\dot\gamma\rangle\right|_{s=0}\\
    &=2\langle\nabla_{\dot c}(tT_{c(t)}),\dot c\rangle
     =2\langle T_{c(t)},\dot c\rangle<0,
\end{align*}
所以, 对于充分小的正数 $s$, $\gamma(s,t)$ 为指向未来的类时曲线. 因此
\[
  q=\lim_{s\to0^+}\gamma(s,1)\in\overline{I^+(p)}.
\]
\item 显然 $I^+(p)\subset\operatorname{int}J^+(p)$. 为了证明 $\operatorname{int}J^+(p)\subset I^+(p)$, 取 $q\in\operatorname{int}J^+(p)$. 取 $q$ 的一个测地凸邻域 $V\subset J^+(p)$, 不难看出, 存在一点 $r\in J^+(p)$, 能沿指向未来的类时曲线连到 $q$. 若 $r_n\in I^+(p)$ 且 $r_n\to r$, 则 $n$ 充分大时有 $r_n\in I^-(q)$ (因而 $q\in I^+(r_n)$). 于是 $q\in I^+(p)$.
\item 显然 $I^+(r)$ 是 $J^+(p)$ 的开子集. 因此 $I^+(r)$ 是 $\operatorname{int}J^+(p)=I^+(p)$ 的开子集.
\item 这等价于证明: 若 $r\in J^-(q)$ 且 $p\in I^-(r)$, 则 $p\in I^-(q)$. 证明同上.
\item 若 $M$ 是二维 Minkowski 时空对由 $f(t,x)=(t+1,x)$ 生成的等距变换群取商所得的空间, 则对任意点 $p\in M$, 有 $I^+(p)=M$.
\end{enumerate}

\RGsolution{ex:small-perturbations-of-stably-causal-spacetime}{small-perturbations-of-stably-causal-spacetime}
对于每一点 $p\in K:=\operatorname{svpp}h$, 在它周围选取局部坐标. 由连续性, 在 $p$ 的某个邻域 $V_p$ 内, 当 $\varepsilon\in(-\varepsilon_p,\varepsilon_p)$ 时, $g_\varepsilon$ 在这些坐标中的分量矩阵 $(g_{\mu\nu}+\varepsilon h_{\mu\nu})$ 有一个负特征值和三个正特征值. 由于 $\{V_p\}_{p\in K}$ 是 $K$ 的开覆盖, 可取有限子覆盖 $\{V_{p_1},\ldots,V_{p_N}\}$. 令 $\varepsilon_0=\min\{\varepsilon_{p_1},\ldots,\varepsilon_{p_N}\}$, 显然当 $|\varepsilon|<\varepsilon_0$ 时, $g_\varepsilon$ 为 Lorentz 度量.

设 $t:M\to\mathbb{R}$ 为 $g$ 的全局时间函数. 由于 $K$ 紧, 在 $K$ 上有
\[
  g(\operatorname{grad}t,\operatorname{grad}t)\leq-\delta
\]
对某个 $\delta>0$ 成立. 考虑由
\[
  f(\varepsilon,p)=g_\varepsilon\bigl((\operatorname{grad}_\varepsilon t)_p,
                                    (\operatorname{grad}_\varepsilon t)_p\bigr)
\]
定义的映射 $f:(-\varepsilon_0,\varepsilon_0)\times K\to\mathbb{R}$, 其中 $\operatorname{grad}_\varepsilon$ 是关于度量 $g_\varepsilon$ 的梯度. 显然, 该映射连续, 且对所有 $p\in K$ 满足
\[
  f(0,p)\leq-\delta.
\]
因此存在 $\varepsilon_1\in(o,\varepsilon_0)$, 使得当 $|\varepsilon|<\varepsilon_1$ 且 $p\in K$ 时,
\[
  f(\varepsilon,p)<0.
\]
因此, 若 $|\varepsilon|<\varepsilon_1$, 则 $t$ 仍然是 $g_\varepsilon$ 的全局时间函数, 所以 $(M,g_\varepsilon)$ 满足时序条件.

\RGsolution{ex:future-domain-of-dependence-examples}{future-domain-of-dependence-examples}
\begin{enumerate}[label=(\alph*),ref=\alph*,leftmargin=*]
\item 显然.
\item 例如, 在二维 Minkowski 时空中取 $S=\{0\}\times[-1,1]$, 则 $D^+(S)$ 是顶点为 $(0,-1)$, $(0,1)$ 及 $(1,0)$ 的闭三角形, 因而不是开集.
\item 例如, 在二维 Minkowski 时空中取 $S=\{0\}\times(-1,1)$, 则 $D^+(S)$ 是 $S$ 与顶点为 $(0,-1)$, $(0,1)$ 及 $(1,0)$ 的开三角形之并, 因而不是闭集.
\end{enumerate}

\RGsolution{ex:global-hyperbolicity-of-product-spacetime}{global-hyperbolicity-of-product-spacetime}
先假定 $(\Sigma,h)$ 完备. 设 $c:I\to M$ 是一条不可延拓的因果曲线. 因为 $t$ 显然是时间函数, 不妨假定 $c$ 以 $t$ 为参数. 因此 $c(t)=(t,\gamma(t))$, 其中 $\gamma:I\to\Sigma$. $c$ 为类时曲线等价于
\[
  h(\dot\gamma,\dot\gamma)\leq1.
\]
设 $t_n\in I$ 是一个有界递增序列, 极限为 $t_0$. 若 $d$ 为 $\Sigma$ 上的 Riemann 距离且 $n>m$, 则
\[
  d(\gamma(t_n),\gamma(t_m))=\int_{t_m}^{t_n}(h(\dot\gamma,\dot\gamma))^{\frac{1}{2}}\,\mathrm{d}t
     \leq t_n-t_m.
\]
所以 $\{\gamma(t_n)\}$ 是 Cauchy 序列, 必收敛, 从而 $\lim_{t\to t_0}\gamma(t)$ 存在. 由于 $c$ 不可延拓, 可知 $I$ 无上界. 同理 $I$ 无下界, 所以 $I=\mathbb{R}$. 可见每条不可延拓的因果曲线都与 $t$ 的每个水平集相交, 因此 $(M,g)$ 整体双曲.

再假定 $(\Sigma,h)$ 不完备, 但 $(M,g)$ 整体双曲. 设 $\gamma:I\to\Sigma$ 为以弧长参数化且不能延拓到 $t\geq0$ 的测地线. 则 $c(t)=(t_0+t,\gamma(t))$ 是一条向未来不可延拓的因果曲线, 位于 $t<t_0$ 的区域内. 由于 $\frac{\partial}{\partial t}$ 的积分曲线是类时曲线, 任何 Cauchy 曲面 $S\subset M$ 都是某个光滑函数 $f:\Sigma\to\mathbb{R}$ 的图像
\[
  S=\{(t,p)\in M\mid t=f(p)\}.
\]
$c$ 到 $S$ 上的投影 $d:I\to S$ 由 $d(t)=(f(\gamma(t)),\gamma(t))$ 给出, 必为类空曲线, 因而满足 $\dot f^2<h(\dot\gamma,\dot\gamma)=1$ (其中 $f(t):=f(\gamma(t))$). 这意味着极限 $f_0:=\lim_{t\to0}f(t)$ 存在. 因为 $|\dot f|<1$, 对所有 $t<0$ 都有 $f(t)>f_0+t$. 取 $t_0=f_0-1$, 便能保证 $c$ 与 $d$ 不相交, 即 $c$ 与 $S$ 不相交. 这就产生矛盾, 说明 $(M,g)$ 不可能整体双曲.

\RGsolution{ex:globally-hyperbolic-spacetime-splitting}{globally-hyperbolic-spacetime-splitting}
设 $t:M\to\mathbb{R}$ 为时间函数, 其水平集是 Cauchy 曲面. 因为 $\operatorname{grad}t$ 不为零, $t$ 不能取到最大值或最小值, 所以其值域必为开区间 $I\subset\mathbb{R}$. 若 $f:I\to\mathbb{R}$ 为微分同胚, 则 $f\circ t$ 也是一个水平集为 Cauchy 曲面的时间函数. 因此可假定 $I=\mathbb{R}$. 定义
\[
  X:=\frac{\operatorname{grad}t}{\langle\operatorname{grad}t,\operatorname{grad}t\rangle},
\]
从而
\[
  X\cdot t=\langle X,\operatorname{grad}t\rangle=1.
\]
由于 $X$ 是类时的, 其积分曲线必穿过时间函数 $t$ 的所有水平集 (因为 $(M,g)$ 整体双曲), 所以 $X$ 是完备的. 若 $\psi_t$ 为 $X$ 的流, 则不难验证由 $F(t,p)=\psi_t(p)$ 定义的 $F:\mathbb{R}\times S\to M$ 为微分同胚.
\end{description}
